\chapter{Introducci\'on}

El an\'alisis de im\'agenes m\'edicas constituye una de las \'areas donde la ingenier\'ia inform\'atica ha adquirido un rol cada vez m\'as relevante. En particular, la aplicaci\'on de t\'ecnicas de procesamiento de im\'agenes e inteligencia artificial permite asistir tareas vinculadas con la detecci\'on, segmentaci\'on, medici\'on y visualizaci\'on de estructuras anat\'omicas. Estas herramientas no buscan reemplazar el criterio profesional, sino complementar el trabajo del especialista mediante resultados cuantitativos, trazables y revisables.

Dentro de este contexto, la resonancia magn\'etica de columna lumbar representa una modalidad de estudio ampliamente utilizada para observar estructuras anat\'omicas como cuerpos vertebrales, discos intervertebrales y canal espinal. La revisi\'on de este tipo de im\'agenes requiere experiencia profesional, interpretaci\'on anat\'omica y, en muchos casos, mediciones o comparaciones visuales entre niveles. Sin embargo, parte del an\'alisis estructural contin\'ua dependiendo de observaci\'on manual, delimitaci\'on visual de regiones de inter\'es y mediciones realizadas de forma no completamente automatizada.

Esta situaci\'on presenta una oportunidad de desarrollo desde la ingenier\'ia inform\'atica: construir una herramienta que permita asistir el an\'alisis de resonancias magn\'eticas lumbares sagitales mediante segmentaci\'on autom\'atica de estructuras relevantes, extracci\'on de mediciones geom\'etricas simples, visualizaci\'on superpuesta de resultados y generaci\'on de una salida estructurada editable. El prop\'osito no es emitir diagn\'osticos ni reemplazar al profesional de salud, sino ofrecer un soporte operativo que permita organizar informaci\'on anat\'omica y cuantitativa de manera revisable.

El presente Proyecto Final de Ingenier\'ia propone el dise\~no y desarrollo de un prototipo funcional de software orientado al an\'alisis asistido de resonancias magn\'eticas lumbares sagitales. El sistema se plantea como una herramienta acad\'emica y experimental que trabaja sobre im\'agenes de RM lumbar, segmenta estructuras anat\'omicas previamente definidas, obtiene mediciones cuantitativas acotadas y presenta los resultados en una interfaz que permita su revisi\'on por parte de un profesional.

El proyecto se apoya principalmente en datasets p\'ublicos anotados, en particular SPIDER, que contiene resonancias magn\'eticas lumbares sagitales con m\'ascaras de referencia para v\'ertebras, discos intervertebrales y canal espinal. Esta decisi\'on permite desarrollar y evaluar t\'ecnicamente el prototipo sin depender inicialmente de datos cl\'inicos privados, favoreciendo la reproducibilidad del trabajo y reduciendo riesgos asociados al tratamiento de informaci\'on sensible.

Desde el punto de vista del alcance, el prototipo se limita a una primera versi\'on funcional o MVP. Este MVP contempla la carga y normalizaci\'on de estudios, segmentaci\'on de estructuras anat\'omicas seleccionadas, extracci\'on de mediciones geom\'etricas simples, visualizaci\'on superpuesta de las m\'ascaras sobre la imagen original y generaci\'on de una salida estructurada editable. Quedan fuera del alcance el diagn\'ostico cl\'inico aut\'onomo, la recomendaci\'on terap\'eutica, la validaci\'on cl\'inica multic\'entrica, la integraci\'on hospitalaria completa y la certificaci\'on regulatoria.

La organizaci\'on del documento sigue una progresi\'on desde la definici\'on del problema hacia el desarrollo y validaci\'on del prototipo. En primer lugar, se presenta el planteamiento del problema, la justificaci\'on, los objetivos, el alcance y la descripci\'on general de la soluci\'on propuesta. Luego se desarrollar\'an el estado del arte y el marco te\'orico, que fundamentar\'an t\'ecnica y conceptualmente el proyecto. Posteriormente se incorporar\'an el user research, los requerimientos, la arquitectura, la metodolog\'ia de desarrollo, la implementaci\'on, las pruebas, los resultados y las conclusiones.

De esta manera, el trabajo busca integrar conocimientos de procesamiento de im\'agenes m\'edicas, aprendizaje autom\'atico, dise\~no de software, interacci\'on humano-computadora y validaci\'on t\'ecnica, dentro de un prototipo acotado y orientado a un problema concreto: asistir el an\'alisis estructural de resonancias magn\'eticas lumbares sagitales mediante segmentaci\'on, medici\'on y revisi\'on profesional.
